\documentclass[11pt]{article}
\usepackage[margin=1in]{geometry}
\usepackage{amsmath,amssymb}
\usepackage{xcolor}

\setlength{\parindent}{1.5em}

\newcommand{\approvalbox}{%
\begin{center}
\fbox{\begin{minipage}{0.85\textwidth}\centering
\textbf{Approval:} \quad $\square$ Student approved \qquad\qquad $\square$ Teaching staff approved
\end{minipage}}
\end{center}
\vspace{6pt}
}

\begin{document}

\begin{center}
{\Large Solutions --- Lab 1: Modeling and Identification of the Parrot AR.Drone}\\[6pt]
{\large Questions 2.1 -- 2.9}
\end{center}

\vspace{18pt}


\vspace{18pt}

\section*{Setup and assumptions}

We use the notation of the handout: $\mathbf{f}_T = \begin{bmatrix} f_1 & f_2 & f_3 & f_4 \end{bmatrix}^T$ for the individual rotor thrusts, $n_i = c f_i$ for the individual reaction torques, $r$ for the arm length (distance from the center of mass to each rotor), and $\mathbf{f}_g = mg\mathbf{e}_3$ for gravity expressed in $\{I\}$.


\section*{2.1 -- Expressions for $\mathbf{f}$ and $\mathbf{n}$}

\approvalbox

\noindent\textbf{Translational force.} The total external force in $\{B\}$ is the sum of gravity (expressed in $\{I\}$, rotated into $\{B\}$) and the resultant of the four rotor thrusts (already aligned with $\{B\}$, all along $-z_B$):
\[
\mathbf{f} = R^T(\lambda)\,\mathbf{f}_g + N\mathbf{f}_T.
\]
Hence
\[
M = R^T(\lambda) \in SO(3), \qquad N = -\mathbf{e}_3 \begin{bmatrix} 1 & 1 & 1 & 1 \end{bmatrix} = \begin{bmatrix} 0 & 0 & 0 & 0 \\ 0 & 0 & 0 & 0 \\ -1 & -1 & -1 & -1 \end{bmatrix} \in \mathbb{R}^{3\times4}.
\]
($M$ depends on attitude because gravity is defined in $\{I\}$ and must be rotated into $\{B\}$; $N$ is constant because every thrust is already expressed in $\{B\}$ along the same axis.)

\noindent\textbf{Torque.} Each rotor $i$ produces (a) a moment from its thrust acting at moment arm $r$ about the two horizontal axes, and (b) a reaction torque $n_i = cf_i$ about $z_B$, alternating in sign with the spin direction. With the geometric convention above ($\mathbf{r}_i \times (-f_i\mathbf{e}_3)$ for each rotor position $\mathbf{r}_i$):
\[
n_x = r(f_2 - f_4), \qquad n_y = r(f_1 - f_3), \qquad n_z = c(f_1 - f_2 + f_3 - f_4),
\]
so that
\[
\mathbf{n} = P\mathbf{f}_T, \qquad P = \begin{bmatrix} 0 & r & 0 & -r \\ r & 0 & -r & 0 \\ c & -c & c & -c \end{bmatrix} \in \mathbb{R}^{3\times4}.
\]

\section*{2.2 -- Redefining $\{B\}$ with $x_B$ at $45^\circ$}

\approvalbox

Rotating the body axes by $45^\circ$ about $z_B$ (new $x_B'$ midway between the arms of rotors 1 and 2) does \emph{not} change $\mathbf{f}$: every thrust still acts along $-z_B$ (unchanged, since the rotation is about $z_B$ itself), and $R^T(\lambda)$ is still just the attitude rotation expressed with respect to whichever axes are called $\{B\}$ --- its functional form does not depend on this relabeling. Hence
\[
M,\ N \text{ unchanged.}
\]

$\mathbf{n} = P\mathbf{f}_T$ \emph{does} change, because the horizontal position of every rotor relative to the (now rotated) $x_B, y_B$ axes changes. The new $x_B'$ is midway between the old $+x_B$ (rotor 1's arm) and old $-y_B$ (rotor 2's arm), i.e. $x_B' = (x_B - y_B)/\sqrt{2}$, $y_B' = (x_B + y_B)/\sqrt{2}$. Expressing each rotor position in the new basis gives $\mathbf{r}_1' = \frac{r\sqrt{2}}{2}(1,1,0)$, $\mathbf{r}_2' = \frac{r\sqrt{2}}{2}(1,-1,0)$, $\mathbf{r}_3' = \frac{r\sqrt{2}}{2}(-1,-1,0)$, $\mathbf{r}_4' = \frac{r\sqrt{2}}{2}(-1,1,0)$ --- each rotor sits at $45^\circ$ from both axes, so each thrust now contributes to \emph{both} $n_x$ and $n_y$ (instead of purely one, as before). Using $\mathbf{r} \times (-f\mathbf{e}_3) = (-fr_y, fr_x, 0)$ for each rotor:
\[
n_x = \frac{r\sqrt{2}}{2}(f_2 + f_3 - f_1 - f_4), \qquad n_y = \frac{r\sqrt{2}}{2}(f_1 + f_2 - f_3 - f_4),
\]
while $n_z$ is unchanged (yaw reaction torque depends only on spin direction, not on the axes' in-plane orientation):
\[
n_z = c(f_1 - f_2 + f_3 - f_4).
\]
So
\[
P' = \begin{bmatrix} -\frac{r\sqrt{2}}{2} & \frac{r\sqrt{2}}{2} & \frac{r\sqrt{2}}{2} & -\frac{r\sqrt{2}}{2} \\[6pt] \frac{r\sqrt{2}}{2} & \frac{r\sqrt{2}}{2} & -\frac{r\sqrt{2}}{2} & -\frac{r\sqrt{2}}{2} \\[6pt] c & -c & c & -c \end{bmatrix} \neq P.
\]


\section*{2.3 -- Input transformation $\mathbf{u} = L\mathbf{f}_T$}

\approvalbox

Stacking $T = \sum_i f_i$ (from 2.1, first row all ones) on top of $\mathbf{n} = P\mathbf{f}_T$ (from 2.1) gives
\[
\mathbf{u} = \begin{bmatrix} T \\ \mathbf{n} \end{bmatrix} = L\mathbf{f}_T, \qquad L = \begin{bmatrix} 1 & 1 & 1 & 1 \\ 0 & r & 0 & -r \\ r & 0 & -r & 0 \\ c & -c & c & -c \end{bmatrix}.
\]

\noindent\textbf{Usefulness.} $L$ is a constant, invertible matrix ($\det L \neq 0$ for $r, c \neq 0$), so the map $\mathbf{f}_T \leftrightarrow \mathbf{u}$ is good. Its point is that the rigid-body dynamics (1)--(3) only ever depend on $\mathbf{f}_T$ through exactly these four combinations: $T$ enters the translational dynamics (scaled along $z_B$) and $\mathbf{n}$ enters the rotational dynamics directly. The individual $f_i$ are physically coupled (each one affects both $T$ and one or two components of $\mathbf{n}$); $\mathbf{u}$ decouples the problem into four independent virtual channels --- one for altitude/thrust and three for the body-frame moments (roll, pitch, yaw). This lets the height and attitude controllers be designed directly in terms of $T$ and $\mathbf{n}$ (as done throughout the rest of the guide), with the physical per-rotor commands recovered only at the very last step via $L^{-1}$ --- i.e., $L$ is exactly the control-allocation / mixer stage of a real flight controller. The only caveat is that not every $\mathbf{u}$ is achievable, since $\mathbf{f}_T$ must remain nonnegative and bounded.

\noindent\textbf{Inverse transformation.} Solving $L\mathbf{f}_T = \mathbf{u}$ directly (summing/subtracting pairs of rows):
\[
f_1 = \frac{T}{4} + \frac{n_y}{2r} + \frac{n_z}{4c}, \qquad f_2 = \frac{T}{4} + \frac{n_x}{2r} - \frac{n_z}{4c},
\]
\[
f_3 = \frac{T}{4} - \frac{n_y}{2r} + \frac{n_z}{4c}, \qquad f_4 = \frac{T}{4} - \frac{n_x}{2r} - \frac{n_z}{4c},
\]
i.e. $\mathbf{f}_T = L^{-1}\mathbf{u}$ with
\[
L^{-1} = \begin{bmatrix} 1/4 & 0 & 1/(2r) & 1/(4c) \\ 1/4 & 1/(2r) & 0 & -1/(4c) \\ 1/4 & 0 & -1/(2r) & 1/(4c) \\ 1/4 & -1/(2r) & 0 & -1/(4c) \end{bmatrix}.
\]

Each rotor receives a quarter of the total thrust plus/minus the roll, pitch and yaw contributions needed to produce the requested $\mathbf{n}$ --- the familiar quadrotor mixer formula.

\section*{2.4 -- Nonlinear state-space dynamics $\dot{x} = g(x,u)$}

\approvalbox

\[
f = mg\,R^T(\lambda)\mathbf{e}_3 - T\mathbf{e}_3.
\]

\noindent\textbf{Kinematics for ${}^B\!\dot p$.} Differentiating ${}^B\!p = R^T p$,
\[
{}^B\!\dot p = \dot R^T p + R^T \dot p.
\]
By definition of the body angular velocity (slide 24):  $\dot R = RS(\omega)$. Using $S(\cdot)^T = -S(\cdot)$ (suggestion),
\[
\dot R^T = S(\omega)^T R^T = -S(\omega)R^T.
\]
Substituting, and using $\dot p = Rv$ (given kinematics) so $R^T\dot p = R^TRv = v$:
\[
{}^B\!\dot p = -S(\omega)R^Tp + v = v - S(\omega)\,{}^B\!p.
\]

\noindent\textbf{Assembling $g(x,u)$.} Combining the kinematics ($\dot p = Rv$, $\dot\lambda = Q(\lambda)\omega$) with the dynamics (1) and the substitutions above:
\[
\dot x = g(x,u) =
\begin{bmatrix}
{}^B\!\dot p \\[4pt] \dot v \\[4pt] \dot\lambda \\[4pt] \dot\omega
\end{bmatrix}
=
\begin{bmatrix}
v - S(\omega)\,{}^B\!p \\[4pt]
-S(\omega)v + g\,R^T(\lambda)\mathbf{e}_3 - (T/m)\mathbf{e}_3 \\[4pt]
Q(\lambda)\,\omega \\[4pt]
J^{-1}\big(u_{2:4} - S(\omega)J\omega\big)
\end{bmatrix},
\]

\section*{2.5 -- Equilibrium point for hover}

\approvalbox

\noindent \textbf{Pedro's proposal: }

By definition, $(x_0,u_0)$ is an equilibrium of $\dot x = g(x,u)$ if and only if $g(x_0,u_0)=0$, so the four blocks derived in 2.4 must vanish simultaneously:
\begin{align*}
v_0 - S(\omega_0)\,{}^B\!p_0 &= 0, &
-S(\omega_0)v_0 + gR^T(\lambda_0)\mathbf{e}_3 - \tfrac{T_0}{m}\mathbf{e}_3 &= 0, \\
Q(\lambda_0)\,\omega_0 &= 0, &
n_0 - S(\omega_0)J\omega_0 &= 0.
\end{align*}

We solve these one at a time. Since $Q(\lambda_0)$ is invertible for any sensible hover attitude ($\theta_0\neq\pm\pi/2$), the third equation forces $\omega_0=0$. Plugging this back into the fourth gives $n_0=0$: no net moment at equilibrium, as expected. The first equation then collapses to $v_0=0$.

The second equation is where the actual work happens. With $\omega_0=0$ it reduces to $gR^T(\lambda_0)\mathbf{e}_3 = (T_0/m)\mathbf{e}_3$. Now $R^T(\lambda)\mathbf{e}_3$ is the third row of $R(\lambda)$ read as a column, namely $[-\sin\theta,\ \sin\phi\cos\theta,\ \cos\phi\cos\theta]^T$, so componentwise
\[
-g\sin\theta_0 = 0, \qquad g\sin\phi_0\cos\theta_0 = 0, \qquad g\cos\phi_0\cos\theta_0 = T_0/m.
\]
Restricting to the physical branch $\theta_0\in(-\pi/2,\pi/2)$, the first equation gives $\theta_0=0$, which reduces the second to $\sin\phi_0=0$ and hence $\phi_0=0$, and finally the third gives $T_0=mg$. Note that $\psi_0$ never appears: a rotation about $z_B$ leaves $\mathbf{e}_3$ fixed, so any yaw is admissible and $\lambda_0=[0,\ 0,\ \psi_0]^T$ with $\psi_0$ free.

Collecting everything,
\[
x_0 = \begin{bmatrix} {}^B\!p_0 \\ v_0 \\ \lambda_0 \\ \omega_0\end{bmatrix} = \begin{bmatrix} R_z^T(\psi_0)\,p_0 \\ \mathbf{0}_3 \\ [0,\ 0,\ \psi_0]^T \\ \mathbf{0}_3 \end{bmatrix}, \qquad u_0 = \begin{bmatrix} T_0 \\ n_0\end{bmatrix} = \begin{bmatrix} mg \\ \mathbf{0}_3\end{bmatrix},
\]
using ${}^B\!p_0=R^T(\lambda_0)p_0=R_z^T(\psi_0)p_0$ since $R(\lambda_0)$ reduces to a pure yaw rotation.

Physically this is exactly what one would expect: at hover the drone has no linear or angular velocity, no net torque, and no roll or pitch, so the thrust vector points straight up and cancels gravity, $T_0=mg$. Both the yaw and the position are free because gravity and the thrust direction (along $z_B$) are invariant under rotation about $z_B$ --- any heading, at any point in space, is a valid hover. Feeding $u_0$ back through the mixer from 2.3 gives $\mathbf{f}_{T0}=L^{-1}u_0=\tfrac{mg}{4}\mathbf{1}_4$, i.e.\ the four rotors share the load equally, as they should for a symmetric quadrotor holding position without yawing.

\section*{2.6 -- Linearized dynamics $\dot{\tilde x}=A\tilde x+B\tilde u$}

\approvalbox



\noindent \textbf{Pedro's proposal: }


\noindent\textbf{Linearization strategy.} Substitute $x=x_0+\tilde x$, $u=u_0+\tilde u$ into each block of $g$ and keep only terms that are zeroth order (which cancel, since $g(x_0,u_0)=0$) or first order in $(\tilde x,\tilde u)$; any product of two or more tilde quantities is second order and is dropped. Since $v_0=\omega_0=0$ and $\tilde p_0=0$ (from 2.5), most quadratic/bilinear terms below vanish identically -- this is exactly the Jacobian $\partial g/\partial x,\partial g/\partial u$ evaluated at $(x_0,u_0)$, just obtained by direct substitution instead of symbolic differentiation.

\noindent\textbf{Assembling $A$ and $B$} (state order $\tilde x=[\tilde p;\tilde v;\tilde\lambda;\tilde\omega]$, input order $\tilde u=[\tilde T;\tilde n]$):
\[
A=\begin{bmatrix}
0_3 & I_3 & 0_3 & 0_3\\
0_3 & 0_3 & A_{v\lambda} & 0_3\\
0_3 & 0_3 & 0_3 & I_3\\
0_3 & 0_3 & 0_3 & 0_3
\end{bmatrix},\qquad
A_{v\lambda}=\begin{bmatrix}0 & -g & 0\\ g & 0 & 0\\ 0 & 0 & 0\end{bmatrix},
\]
\[
B=\begin{bmatrix}
0_{3\times1} & 0_{3\times3}\\
b_T & 0_{3\times3}\\
0_{3\times1} & 0_{3\times3}\\
0_{3\times1} & J^{-1}
\end{bmatrix},\qquad b_T=\begin{bmatrix}0\\0\\-1/m\end{bmatrix}.
\]

\section*{2.7 -- Transfer functions from the linearized dynamics}

\approvalbox

With $J$ diagonal, each block of $A,B$ from 2.6 decouples into three independent second-order chains, read off row by row.

\noindent\textbf{Attitude channels.} $\dot{\tilde\lambda}=\tilde\omega$ and $\dot{\tilde\omega}=J^{-1}\tilde n$ give, componentwise, $\ddot\phi=\tilde n_x/J_{xx}$, $\ddot\theta=\tilde n_y/J_{yy}$, $\ddot\psi=\tilde n_z/J_{zz}$. Hence
\[
G_\phi(s)=\frac{1}{J_{xx}s^2}, \qquad G_\theta(s)=\frac{1}{J_{yy}s^2}, \qquad G_\psi(s)=\frac{1}{J_{zz}s^2}.
\]

\noindent\textbf{Vertical channel.} $\dot{\tilde z}=\tilde w$ and $\dot{\tilde w}=-\tilde T/m$ (third rows of the $\dot{\tilde p}=\tilde v$ and $\dot{\tilde v}=A_{v\lambda}\tilde\lambda+b_T\tilde T$ blocks) give $\ddot{\tilde z}=-\tilde T/m$, so
\[
G_z(s)=\frac{Z(s)}{T(s)}=-\frac{1}{ms^2}.
\]

\noindent\textbf{Horizontal channels.} From $A_{v\lambda}$: $\dot{\tilde u}=-g\tilde\theta$ and $\dot{\tilde v}=g\tilde\phi$ (first two rows), together with $\dot{\tilde x}=\tilde u$, $\dot{\tilde y}=\tilde v$, give $\ddot{\tilde x}=-g\tilde\theta$ and $\ddot{\tilde y}=g\tilde\phi$. Substituting $\Theta(s)=G_\theta(s)N_y(s)$ and $\Phi(s)=G_\phi(s)N_x(s)$:
\[
G_x(s)=\frac{X(s)}{N_y(s)}=-\frac{g}{J_{yy}s^4}, \qquad G_y(s)=\frac{Y(s)}{N_x(s)}=\frac{g}{J_{xx}s^4}.
\]

Note $\tilde x,\tilde y,\tilde z$ here are the components of ${}^B\tilde p$ (the state chosen in 2.4--2.6), not of the inertial $p$ directly -- relevant for the discussion in 2.8.

\section*{2.8 -- Relating $G_\phi,G_y$ and $G_\theta,G_x$}

\approvalbox

Dividing the horizontal transfer functions from 2.7 by the corresponding attitude ones:
\[
G_y(s)=\frac{g}{s^2}\,G_\phi(s), \qquad G_x(s)=-\frac{g}{s^2}\,G_\theta(s).
\]

\noindent\textbf{Physical interpretation.} Tilting about $x_B$ (roll, $\phi$) tips the thrust vector toward $y_B$, producing a horizontal acceleration $g\phi$; tilting about $y_B$ (pitch, $\theta$) tips it toward $-x_B$, producing $-g\theta$. Position is acceleration integrated twice, so each horizontal transfer function is exactly its attitude counterpart cascaded with two extra integrators and a gain $g$: attitude must first build up before it can move the vehicle, it never happens instantaneously.

\noindent\textbf{Effect of $\psi_0$.} $\tilde x,\tilde y$ are body-frame quantities (${}^B\tilde p$), so they relate to the actual inertial horizontal displacement via $R_z(\psi_0)$, which mixes $x$ and $y$ whenever $\psi_0\neq0$ (unlike $\tilde z$, which is unaffected by yaw, cf. 2.9). A roll input therefore does not move the vehicle purely along inertial North or East once $\psi_0\neq0$: the decoupled pairing $G_\phi\!\leftrightarrow\!G_y$, $G_\theta\!\leftrightarrow\!G_x$ only holds in body coordinates, and any outer position-loop working in the inertial frame must rotate its commanded accelerations by $\psi_0$ before splitting them into roll/pitch references -- exactly the yaw-compensation stage present between the position and attitude loops of a real flight controller.

\section*{2.9 -- Height control schemes}

\approvalbox

\noindent\textbf{Plant.} With $h=-z$ and $h_0=-z_0$, $\tilde h=-\tilde z$; since equilibrium hover has $\phi_0=\theta_0=0$, yaw alone does not affect the $z$-component (cf. 2.8), so $\tilde z={}^B\tilde p_z$ exactly and $H(s)=-\tilde Z(s)$. From 2.7, $G_z(s)=-1/(ms^2)$, so
\[
\frac{H(s)}{T(s)}=-G_z(s)=\frac{1}{ms^2},
\]
matching the cascade $\frac{1}{ms}\!\to\!\frac1s$ (force$\to\dot h\to h$) shown in Figs.\ 2--3.

\noindent\textbf{Closed-loop transfer functions.} For Fig.\ 3 (series PD), $T=(k_p+k_ds)(H_{ref}-H)$ directly gives
\[
\frac{H(s)}{H_{ref}(s)}=\frac{k_ds+k_p}{ms^2+k_ds+k_p}.
\]
For Fig.\ 2 (cascade), the inner loop around $1/(ms)$ with feedback $k_d$ gives $\dot H=\frac{1}{ms}\big[k_p(H_{ref}-H)-k_ds H\big]$; solving for $H$:
\[
\frac{H(s)}{H_{ref}(s)}=\frac{k_p}{ms^2+k_ds+k_p}.
\]
The two share the same denominator (same poles, same root-locus in $k_p$ for fixed $k_d$, cf.\ 2.10) but Fig.\ 3 has an extra zero at $s=-k_p/k_d$.

\noindent\textbf{Differences.} Writing $G_3=G_2+\frac{k_ds}{ms^2+k_ds+k_p}$, the Fig.\ 3 step response equals the Fig.\ 2 response plus a term that starts at $0$ with slope $k_dA/m>0$: Fig.\ 3 always overshoots more for the same $(k_p,k_d)$. Computing $T(s)$ for a step $H_{ref}(s)=A/s$ in each scheme and applying the initial value theorem gives $T(0^+)=k_pA$ (finite) for Fig.\ 2, but $T(0^+)\to\infty$ for Fig.\ 3 -- the series PD differentiates the reference itself, so a step $h_{ref}$ demands an infinite instantaneous thrust ("derivative kick"), while the cascade never differentiates $h_{ref}$, only the measured $h$. Likewise, if $\dot h$ in Fig.\ 2 comes from a genuine rate measurement (rather than from numerically differentiating $h$), its noise-to-actuator gain is the flat $|k_d|$, whereas Fig.\ 3's $k_ds$ branch has gain $|k_d\omega|$, unbounded in frequency -- it amplifies high-frequency sensor noise that the cascade avoids.

\noindent\textbf{Advantages/disadvantages.} Fig.\ 2 (cascade): smoother, more predictable step response (no zero), no derivative kick, and -- if a rate sensor is available -- no derivative-induced noise amplification; but it needs access to a real $\dot h$ signal and is structurally more complex (nested loop). Fig.\ 3 (series PD): simpler to implement (single compensator) and reacts faster to reference changes, at the cost of more overshoot, an unrealizable derivative kick on step references, and noise amplification if $\dot h$ must be obtained by differentiating a noisy $h$.

\end{document}
